\scanpage{83}
% archive foliation: 82
% Printed face of the reused sheet; image rotated 180 degrees for reading.
\begin{center}$-\ 122\ -$\end{center}
\begin{flushright}$n^{\circ}\ 316$\end{flushright}
$V_{(L)}$ est muni canoniquement d'une structure d'espace vectoriel topologique
métrisable et complet (\uline{Esp.} \uline{vect.} \uline{top.} chap. I, § 2, $n^{\circ}3$, th.2).

Soit $(x_p)_{p\in\mathbf Z}$ une famille d'éléments de $V$ tels que $x_p=0$
pour $p$ inférieur à un certain entier rationnel.
Alors la série $\sum_{-\infty}^{+\infty}x_pT^p$ converge dans $V_{(L)}$,
et tout élément de $V_{(L)}$ se met sous cette forme de manière unique.
(Prendre une base dans $V$).

b) L'espace vectoriel $\mathcal L_L(V_{(L)})=(\mathcal L_K(V))_{(L)}$
est muni canoniquement d'une structure d'espace vectoriel topologique métrisable et complet.
Tout endomorphisme de $V_{(L)}$ se met d'une manière unique sous la forme
$\sum_{-\infty}^{+\infty}u_pT^p$, où $u_p\in\mathcal L_K(V)$,
et où les $u_p$ sont nuls pour $p$ inférieur à un certain entier rationnel.

c) Si $K$ est de caractéristique 0, et si $u\in\mathcal L_K(V)$, on pose
$e^{Tu}=1+\frac{1}{1!}uT+\frac{1}{2!}u^2T^2+\cdots\in\mathcal L_L(V_{(L)})$.
Si $x\in V$ est tel que $ux=0$, on a $e^{Tu}.x=x$.

d) Si $W$ est un autre espace vectoriel de dimension finie sur $K$,
et si $v\in\mathcal L_K(W)$, on a
$e^{T(u\otimes1+1\otimes v)}=e^{Tu}\otimes e^{Tv}$.

e) Déduire de c) et d) que, si $V$ est muni d'une structure d'algèbre
